% Core control FSM. Requires tikz_styles.tex.
\begin{figure}[p]
\centering
\begin{tikzpicture}[x=1cm,y=1cm]
  % Main pipeline
  \node[ctl, minimum width=15mm] (boot) at (0.9,7.6) {BOOT};
  \node[ctl, minimum width=16mm] (fetch) at (3.0,7.6) {FETCH};
  \node[ctl, minimum width=17mm] (dec) at (5.2,7.6) {DECODE};
  \node[ctl, minimum width=15mm] (exec) at (7.3,7.6) {EXEC};
  \node[ctl, minimum width=14mm] (mems) at (9.2,7.6) {MEM};
  \node[ctl, minimum width=14mm] (wb) at (11.1,7.6) {WB};

  % Departures from EXEC / WB
  \node[acc, minimum width=22mm] (cont) at (3.0,5.9) {CONTAINER};
  \node[acc, minimum width=18mm] (sa) at (5.6,5.9) {STRACC};
  \node[ctl, minimum width=14mm] (call) at (8.0,5.9) {CALL};
  \node[ctl, minimum width=18mm] (ret) at (10.6,5.9) {RETURN};

  \node[rf, minimum width=18mm] (spill) at (8.0,4.3) {RF\_SPILL};
  \node[rf, minimum width=16mm] (fill) at (10.6,4.3) {RF\_FILL};
  \node[rf, minimum width=15mm] (init) at (12.7,4.3) {RF\_INIT};
  \node[mem, minimum width=20mm] (cw) at (15.0,4.3) {CODE\_WRITE};

  % Trap and GC, kept on their own rows
  \node[ex, minimum width=20mm] (mar) at (1.2,2.6) {MARSHAL};
  \node[ex, minimum width=16mm] (wait) at (3.8,2.6) {WAIT};
  \node[trap, minimum width=14mm] (halt) at (1.2,1.0) {HALT};

  \node[acc, minimum width=20mm] (ent) at (6.4,2.6) {GC\_ENTER};
  \node[acc, minimum width=20mm] (roots) at (9.0,2.6) {GC\_ROOTS};
  \node[acc, minimum width=16mm] (run) at (11.4,2.6) {GC\_RUN};
  \node[acc, minimum width=18mm] (alloc) at (13.8,2.6) {GC\_ALLOC};

  \draw[arr] (boot) -- (fetch);
  \draw[arr] (fetch) -- (dec);
  \draw[arr] (dec) -- (exec);
  \draw[arr] (exec) -- (mems);
  \draw[arr] (mems) -- (wb);
  \draw[arr] (wb.north) -- ++(0,0.4) -| (fetch.north);

  \draw[arr] (exec.south) -- ++(0,-0.28) -| (sa.north);
  \draw[arr] (dec.south) -- (cont.north);
  \draw[arr] (wb.south) -- ++(0,-0.28) -| (ret.north);
  \draw[arr] (mems.south) -- ++(0,-0.28) -| (call.north);

  \draw[arr] (call) -- (spill);
  \draw[arr] (ret) -- (fill);

  \draw[arr] (cont.south) -- ++(0,-0.45) -| (mar.north);
  \draw[arr] (mar) -- (wait);
  \draw[arr] (wait.south) -- ++(0,-0.35) -| (halt.east);
  \draw[arr] (ent) -- (roots);
  \draw[arr] (roots) -- (run);
  \draw[arr] (run) -- (alloc);

  \node[anno, anchor=south] at (7.0,8.25) {next instruction};
  \node[anno, anchor=west] at (2.15,1.0) {fatal};
  \node[anno, anchor=north] at (10.2,1.85) {empty run list forces a collection};
\end{tikzpicture}
\caption{Control states in \texttt{pycore\_core}.
The top row is the steady path; the arrow above it is the return to fetch.
Solid edges are completions, not the self-loops that wait on \texttt{ack} or \texttt{stall\_o}.
Any latched fatal trap forces \texttt{S\_HALT}.
\texttt{S\_FETCH} holds until \texttt{pycore\_fetch} presents a real instruction
(\texttt{CACHE} skipped, \texttt{EXTENDED\_ARG} folded inside the fetch unit).
Recoverable traps leave the container FSM for \texttt{S\_TRAP\_MARSHAL} when \texttt{EXCORE\_EN=1}.
A collection that starts because an allocation does not fit enters \texttt{S\_GC\_ALLOC} directly.
An empty free-run list continues into \texttt{S\_GC\_ENTER}.}
\label{fig:fsm}
\end{figure}

\begin{figure}[htbp]
\centering
\begin{tikzpicture}[x=0.95cm,y=0.72cm]
  \draw[pyedge] (0,0) -- (14.2,0);
  \foreach \x/\lab in {0/FETCH, 1.7/DEC, 3.4/EXEC, 6.6/MEM, 8.6/WB, 11.2/FETCH} {
    \draw[pyedge] (\x,-0.12) -- (\x,0.16);
    \node[anno, anchor=north] at (\x,-0.16) {\lab};
  }
  \fill[pyctl] (0.15,0.45) rectangle (1.45,1.15);
  \fill[pyctl] (1.85,0.45) rectangle (3.15,1.15);
  \fill[pyacc] (3.55,0.45) rectangle (6.35,1.15);
  \fill[pymem] (6.75,0.45) rectangle (8.35,1.15);
  \fill[pyrf] (8.75,0.45) rectangle (10.95,1.15);
  \node[font=\tiny\sffamily] at (0.8,0.8) {latch};
  \node[font=\tiny\sffamily] at (2.5,0.8) {RF read};
  \node[font=\tiny\sffamily] at (4.95,0.8) {unit busy};
  \node[font=\tiny\sffamily] at (7.55,0.8) {dmem};
  \node[font=\tiny\sffamily] at (9.85,0.8) {RF write};
  \node[anno, anchor=west] at (0,1.45)
    {The next fetch does not start until writeback retires.};
  \node[anno, anchor=west] at (0,-1.25)
    {Multi-cycle arithmetic stretches EXEC. A container or STRACC op replaces MEM and WB.};
\end{tikzpicture}
\caption{One instruction in flight.
There is no EX/MEM or MEM/WB forwarding, no load-use interlock, and no branch flush of younger instructions.
Taken branches redirect \texttt{pycore\_fetch} on the next entry to \texttt{S\_FETCH}.
Decode reads \texttt{tos\_r} combinationally, so the address it forms already includes the previous instruction's stack update.}
\label{fig:inflight}
\end{figure}
